\documentclass[english, parskip=full]{scrartcl}
\usepackage[utf8]{inputenc}  % Kodierung der Datei
\usepackage[english]{babel}  % Sprache auf Deutsch 
\usepackage{color}
\usepackage{graphicx, gensymb}
\usepackage{amsfonts, cancel, amssymb}
\usepackage{amsmath, amsthm, array, esint}
\usepackage{booktabs, multirow}  % schönere Tabellen
\usepackage{caption}  % Anpassung der Captions
\usepackage{scrlayer-scrpage}    % Manipulation des Seitenstils
\pagestyle{scrheadings}  % Stil für die Seite setzen
\automark{section}  % Abschnittsnamen als Seitenbeschriftung 
\setheadsepline{.5pt}    % Kopzeile durch Linie abtrennen
\usepackage[hidelinks]{hyperref}  % Links und weitere PDF-Features
\usepackage{typearea, tikz, pgffor, verbatim}
\usetikzlibrary{calc, arrows.meta,arrows, graphs, quotes, positioning}
\usepackage[cache=false, outputdir=build]{minted}
\usemintedstyle{vs}
\usepackage{placeins} % provides \FloatBarrier

\definecolor{bg}{rgb}{0.9,0.9,0.9}
\newcommand\numberthis{\addtocounter{equation}{1}\tag{\theequation}}
\usepackage{svg}
\usepackage{siunitx}
\usepackage{physics}
\usepackage{tensor}
\renewcommand{\bra}{}
\newcommand{\kom}{}
\newcommand{\brak}{}
\renewcommand{\braket}{}
\newcommand{\intx}{}
\newcommand{\intr}{}
\newcommand{\kla}{}
\newcommand{\klam}{}
\newcommand{\klammer}{}
\newcommand{\oper}{}
\newcommand{\tecxt}{}
\newcommand{\Bra}{}
\newcommand{\Ket}{}
\renewcommand{\i}{\text{i}}
\newcommand{\e}{\text{e}}
\newcommand*{\Fourier}{\textsc{Fourier}\ }
\newcommand*{\F}{\mathcal{F}}
\newcommand*{\sinc}{\text{sinc\,}}
\newcommand{\conv}{\mathop{\scalebox{3.5}{\raisebox{-0.38ex}{\makebox[0.5\width][c]{$\ast$}}}}\limits}

\newcommand*{\titel}{Hamiltonian Ratchet}
\newcommand*{\autor}{Joachim Schwardt}

\hypersetup{pdfauthor={\autor}, pdftitle={\titel}}

%\titlehead{titlehead}
\subject{Stochastic processes}
\title{\titel}
\author{\autor}
\date{\today}

\begin{document}

%\begin{titlepage}
\maketitle  % Titel setzen
%\tableofcontents  % Inhaltsverzeichnis setzen
%\end{titlepage}

\section{Definition}
We have a potential of the form 
\begin{align}
V(x) &= V_0 \klam\left[ \sin\kla\left( \frac{2\pi x}{L} \right) + a\cos\kla\left( \frac{4\pi x}{L} \right) \right],
\end{align}
and execute the equation
\begin{align}
x_{n+1} &= x_n - V'(x_n) \Delta t + \sqrt{2D\Delta t} \xi_n
\end{align}
for many iterations of $n$.
Here $\xi_n$ is ``white noise'' and $D$ is the diffusion constant.
Note that we run for $N_\tau$ periods of length $\tau$, and we discretize each $\tau$ with $S$ steps.
This means that $\Delta t = \frac{1}{S}$.
%Now let us transform to new units, in which $\tau' := 1$. 
%Then we have to shift $S' = \tau S$ and $\Delta t' = \frac{1}{\tau} \Delta t$.
%Let us also fix $L' := 1$.
%Then
%\begin{align}
%V(x') &= V_0 \klam\left[ \sin (2\pi x') + a\cos (2\pi x') \right]
%\end{align}
%and the derivative becomes $\partial_{x'} = L\partial_x$.
%Thus the iteration now reads
%\begin{align}
%x'_{n+1} &= x'_n - L V'(x_n') \Delta t' + \sqrt{\frac{2D\tau\Delta t'}{L^2}} \xi_n.
%\end{align}
%This allows us to identify $V_0' = L V_0$ and $D' = \frac{\tau}{L^2} D$.
%FIXME: can we not just say that $\partial_{x'}= L\partial_x$ and $D' = \frac{\tau}{L^2} D$ follow trivially from the dimensional analysis?
%Also, $\Delta t$ is really an invariant of this transformation.
%This is because we adjust the number of steps in order to keep it constant on purpose, so the arguments above fail.
Now let us transform to new units, in which $\tau' := 1$. 
Then we have to shift $S' = \tau S$ and $\Delta t' = \frac{1}{\tau} \Delta t$.
This means, that we have to adjust $V_0$ and $D$ slightly differently.
Let us also fix $L' := 1$.
Then
\begin{align}
V(x') &= V_0 \klam\left[ \sin (2\pi x') + a\cos (2\pi x') \right]
\end{align}
and the derivative becomes $\partial_{x'} = L\partial_x$.
Diving through by $L$, the iteration now reads
\begin{align}
x'_{n+1} &= x'_n - \tau V'(x_n') \Delta t' + \sqrt{\frac{2D\tau\Delta t'}{L^2}} \xi_n.
\end{align}
This allows us to identify $V_0' = \tau V_0$ and $D' = \frac{\tau}{L^2} D$.

%\begin{thebibliography}{9}
%\bibitem{example} description, \url{https://www.exmapleurl}, day.\,month~2021.
%\end{thebibliography}

\end{document}